% paper/sections/11_benchmarks.tex
%
% §11  多相流ベンチマークと物理現象の検証（Validation）
%
%   §11.1: Balanced-Force と寄生流れ — 静止液滴テスト
%   §11.2: 毛細管波の周期と減衰 — Prosperetti 解析解との比較
%   §11.3: 気泡上昇ベンチマーク（Hysing et al. 2009）
%   §11.4: Rayleigh--Taylor 不安定性（線形成長率検証）

\section{多相流ベンチマークと物理現象の検証}
\label{sec:validation}
\label{sec:benchmarks}

第~\ref{sec:component_verification}章のコンポーネント検証で各部品の正確性を確認したのち，
本章では複数の物理機構が連成する多相流問題に対するフル NS ソルバーの総合的な検証を行う．
ベンチマークは難易度の低い順に配置されており，各ステップで合格してから次のテストに進むことを推奨する．

% ═══════════════════════════════════════════════════════════════════════════════
\subsection{Balanced-Force と寄生流れの抑制：静止液滴テスト}
\label{sec:val_static_drop}
\label{sec:bench_droplet}
% ═══════════════════════════════════════════════════════════════════════════════

\textbf{目的：}
Balanced-Force 離散化（§\ref{sec:balanced_force}）の下で，表面張力と圧力勾配が格子レベルで
釣り合い，数値起源の寄生流れ（parasitic currents）が CSF モデル誤差 $O(\varepsilon^2) \approx O(h^2)$
水準まで抑制されることを定量的に示す．
これは本手法の主要な技術的主張であり，高次精度 CCD の採用効果が直接現れる．

\textbf{設定：}
\begin{itemize}
  \item 正方形領域 $[0,1]^2$，中心 $(0.5, 0.5)$ に半径 $R=0.25$ の静止液滴
  \item 初期速度 $\bm{u} = \bm{0}$；密度比 $\rho_l/\rho_g = 1000$，$We = 1$，$Fr = \infty$（重力なし）
  \item 格子解像度：$N = 32, 64, 128$
\end{itemize}

\textbf{検証指標：}
\begin{enumerate}
  \item \textbf{ラプラス圧力：}
        内外圧差 $|\Delta p| = |p_{\mathrm{in}} - p_{\mathrm{out}}|$ が $\sigma / R$ に一致すること
        （相対誤差 $< 1\%$）．
        符号については §\ref{sec:governing} の規約 $[p] = \sigma\kappa = -\sigma/R < 0$ を参照．
  \item \textbf{寄生流れ収束：}
        $\|\bm{u}_{\mathrm{para}}\|_\infty$ が $O(h^2)$ で収束すること．
        $N = 64$，$\varepsilon \approx 1.5h$ での目標：
        $\|\bm{u}\|_\infty / (\sigma/\mu_l) \lesssim 10^{-5}$（Hysing et al.~\cite{Hysing2009} 水準）．
\end{enumerate}

\begin{tcolorbox}[resultbox, title={寄生流れ振幅の理論的予測（Balanced-Force 条件下）}]
Balanced-Force 条件下（§\ref{sec:balanced_force}）では，圧力勾配 $\bnabla p$ と
表面張力 $\sigma\kappa\bnabla\psi$ に\textbf{同一の CCD 演算子}を用いることで，
離散化誤差 $\mathcal{O}(h^6)$ が両項で相殺される．
残差は CSF モデル誤差のみとなり，有効寄生流れ振幅は
\[
  \|\bm{u}_{\mathrm{para}}\| \sim \frac{\sigma}{\mu_l} C_{\mathrm{model}} \varepsilon^2 = O(h^2)
\]
$h = 1/64$，$\varepsilon \approx 1.5h$ では $\|\bm{u}_{\mathrm{para}}\|/(\sigma/\mu_l) \sim 10^{-5}$ が目標値となる
（非 Balanced-Force 2次手法と比べ約1桁の低減；§\ref{sec:bench_droplet}の詳細解析参照）．
\end{tcolorbox}

\textbf{期待される結果：}
\begin{center}
\renewcommand{\arraystretch}{1.3}
\begin{tabular}{ccc}
\toprule
$N$ & $\|\bm{u}_{\mathrm{para}}\|_\infty / (\sigma/\mu_l)$（目標） & 収束次数 \\
\midrule
$32$  & $\sim 10^{-3}$ & --- \\
$64$  & $\sim 2.5\times10^{-4}$ & $\approx 2$ \\
$128$ & $\sim 6.3\times10^{-5}$ & $\approx 2$ \\
\bottomrule
\end{tabular}
\end{center}

% ═══════════════════════════════════════════════════════════════════════════════
\subsection{毛細管波の周期と減衰：Prosperetti 解析解との比較}
\label{sec:val_capwave}
% ═══════════════════════════════════════════════════════════════════════════════

\textbf{目的：}
粘性・表面張力・圧力の動的連成が正しく実装されていることを，
毛細管波（capillary wave）の振動周期と減衰率に関する
Prosperetti (1981) \cite{Prosperetti1981} の線形解析解との定量比較によって確認する．
解析解が存在する問題であるため，数値誤差を厳密に評価できる．

\textbf{設定：}
\begin{itemize}
  \item 計算領域 $[0, \lambda] \times [0, 2\lambda]$（$\lambda$：波長，基準長さ $\lambda = 1$），
        周期 BC（$x$ 方向），固定壁 BC（$y$ 方向）
  \item 水平界面 $y = \lambda$（領域中央）に微小正弦波擾乱
        $\eta(x,0) = \eta_0 \sin(2\pi x / \lambda)$（$\eta_0/\lambda = 10^{-2}$）
  \item 初期速度 $\bm{u} = \bm{0}$；密度比 $\rho_l/\rho_g = 2$，$\mu_l = \mu_g = \mu$（等粘性）
  \item 検査パラメータ：Laplace 数 $La = \rho_l \sigma \lambda / \mu_l^2$（$La = 10$：過減衰域；$La = 3000$：振動域）
  \item 格子：$N_x = 64$（$\lambda / h = 64$），$N_y = 128$
\end{itemize}

\textbf{Prosperetti 解析解の概要：}
線形化した2流体 NS 方程式に正弦波解 $\eta \propto e^{st}$ を仮定すると，
成長率 $s$ は4次の分散方程式の根として得られる \cite{Prosperetti1981}：
\begin{equation}
  \left(s + 2\nu k^2\right)^2 \!+ \frac{\sigma k^3}{\rho_l + \rho_g}
  = 4\nu^2 k^4 \sqrt{1 + \frac{s}{\nu k^2}}
  \label{eq:prosperetti_dispersion}
\end{equation}
ここで $k = 2\pi/\lambda$ は波数，$\nu = \mu/\rho_l$ は動粘性係数である．
$La \gg 1$ の非粘性極限では $s \to \pm i\omega_0$ に振動解が現れ，
$\omega_0 = \sqrt{\sigma k^3 / (\rho_l + \rho_g)}$ （毛細管波角周波数）となる．

\textbf{検証指標：}
\begin{itemize}
  \item 界面変位振幅 $\eta(t)$ の時間発展を計測し，
        式~\eqref{eq:prosperetti_dispersion} の数値根と比較する
  \item 振動周期・減衰率の相対誤差 $< 1\%$
\end{itemize}

\begin{tcolorbox}[mybox, title={参照パラメータ（Harvie \& Fletcher 2000 より）}]
\begin{center}
\renewcommand{\arraystretch}{1.2}
\begin{tabular}{lll}
\hline
ケース & $La$ & 挙動 \\
\hline
過減衰 & $10$ & 振動なし；指数減衰 \\
臨界減衰付近 & $50$ & 1--2 振動後に減衰 \\
振動域 & $3000$ & 長周期振動 + 緩やかな減衰 \\
\hline
\end{tabular}
\end{center}

毛細管時刻 $t_\sigma = \sqrt{(\rho_l+\rho_g)\lambda^3 / (2\pi\sigma)}$ で時間を無次元化し，
$t \in [0, 5 t_\sigma]$ を積分する（$N_y = 128$ では $t_\sigma \approx 1/\omega_0$）．
\end{tcolorbox}

% ═══════════════════════════════════════════════════════════════════════════════
\subsection{気泡上昇ベンチマーク（Hysing et al. 2009）}
\label{sec:val_rising_bubble}
\label{sec:bench_bubble}
% ═══════════════════════════════════════════════════════════════════════════════

\textbf{目的：}
表面張力・重力・粘性・密度差が同時に作用する現実的な問題で，
ソルバー全体の総合的な安定性と精度を Hysing ら \cite{Hysing2009} の
高精度参照解（複数グループの独立検証による「数値的参照解」）と定量比較する．

\textbf{設定（Test Case 1）：}
\begin{itemize}
  \item 計算領域 $[0, 1] \times [0, 2]$，初期気泡：中心 $(0.5, 0.5)$，半径 $R = 0.25$
  \item $\rho_l / \rho_g = 10$，$\mu_l / \mu_g = 10$，$Re = 35$，$We = 10$（$Fr^2 = 1$）
  \item 側面・底面・上面：スリップ BC（$\partial u/\partial n = 0$，$v = 0$）
  \item 格子：$N_x \times N_y = 64 \times 128$（収束確認：$128 \times 256$）
\end{itemize}

\textbf{検証指標：}
\begin{align}
  V_c(t) &= \frac{\int_{\Omega_b} v \, dV}{|\Omega_b|}
  \quad (\text{気泡重心の上昇速度}) \label{eq:bubble_vc_val} \\
  C_{\mathrm{circ}}(t) &= \frac{\text{等面積円の周長}}{\text{実際の気泡周長}}
  \quad (\text{真球度，= 1 で球形}) \label{eq:bubble_circ_val}
\end{align}
$t = 3$ での $V_c$ と $C_{\mathrm{circ}}$ の Hysing 参照解との相対誤差 $< 1\%$ を目標とする．

\begin{tcolorbox}[resultbox, title={Hysing (2009) Test Case 1 の参照値（数値参照解）}]
\textbf{重要：これらの値は解析解ではなく，複数の独立グループによる「数値的参照解」である．}

\begin{center}
\renewcommand{\arraystretch}{1.3}
\begin{tabular}{lcccc}
\toprule
指標 & TP2D & FreeLIFE & MooNMD & 参照値（平均） \\
\midrule
$V_c(t=3)$ & $0.2417$ & $0.2421$ & $0.2414$ & $\approx 0.2417$ \\
$C_{\mathrm{circ}}(t=3)$ & $0.9011$ & $0.9013$ & $0.9015$ & $\approx 0.9013$ \\
$y_c(t=3)$ & $1.0816$ & $1.0810$ & $1.0812$ & $\approx 1.0813$ \\
\bottomrule
\end{tabular}
\end{center}

グループ間の差 $\Delta V_c \approx 0.001$ ($\approx 0.4\%$) が本ベンチマークの不確かさであり，
本ソルバーの結果がこの範囲内に収まることを合格基準とする．

動格子（ALE）補正の省略については付録~\ref{app:grid_ale} 参照；
本パラメータ設定（$Re = 35$）では ALE 省略の影響は軽微である．
\end{tcolorbox}

% ═══════════════════════════════════════════════════════════════════════════════
\subsection{Rayleigh--Taylor 不安定性（線形成長率検証）}
\label{sec:val_rt}
\label{sec:bench_rt}
% ═══════════════════════════════════════════════════════════════════════════════

\textbf{目的：}
密度差・重力による界面変形の長時間シミュレーション能力を検証する．
線形安定解析から導出された理論成長率（式~\eqref{eq:rt_growth}）と
数値解の初期成長率を比較し，
非線形段階での定性的なバブル--スパイク構造の形成を確認する．

\textbf{設定：}
\begin{itemize}
  \item 計算領域 $[0, \lambda] \times [0, 4\lambda]$（$\lambda = 1$），
        周期 BC（$x$），固定壁 BC（$y$）
  \item 上層（重流体 $\rho_l = 3$）が下層（軽流体 $\rho_g = 1$）の上に乗る不安定成層
        （Atwood 数 $At = (\rho_l-\rho_g)/(\rho_l+\rho_g) = 0.5$）
  \item 初期擾乱：$y_0(x) = 2 + A_0\sin(2\pi x)$（$A_0 = 0.05$）；$\mu_l = \mu_g = 0.01$，$\sigma = 0$
  \item 格子：$64 \times 256$
\end{itemize}

\textbf{線形成長率との比較：}
\begin{equation}
  \omega_{\mathrm{RT}} = \sqrt{\frac{\rho_l - \rho_g}{\rho_l + \rho_g} \cdot \frac{2\pi g}{\lambda}}
  = \sqrt{\pi} \approx 1.7725
  \quad [\text{無次元時間}^{-1}]
  \label{eq:rt_growth}
\end{equation}

振幅が線形域 $A(t) < 0.1\lambda$ にある間（$t \lesssim 1.5$）に数値成長率を測定し，
理論値との相対誤差 $< 5\%$ を確認する
（詳細な導出は §\ref{sec:bench_rt}，線形安定解析の概略は Chandrasekhar \cite{Chandrasekhar1961} 参照）．
